\subsection{A normalized geometric trace for the actual quadratic coefficient}

\begin{proposition}[P-F69: an ordinary abelian-surface trace counterpart]
\label{prop:nonstandardhondatrace}
Retain $A,D,F,P,T$ and the actual source eigenvalues
$\alpha_\pm$ from P-F66.  Fix $e\ge1$, set $q=3^e$, and put
\[
\begin{split}
 B&=(1+q+\cdots+q^9)^2(1+q^2)(1+q^{12}),\\
 E&=q^{10}T(q),\qquad K=\mathbb Q(\sqrt{A(q)}),\\
 \beta_\pm&=B\bigl(E\pm q^2(1+q^8)\sqrt{A(q)}\bigr)
                  =2q^{26}\alpha_\pm(q),\\
 Q&=q^{116}.
\end{split}                                                   \tag{P-F69a}
\]
There exists a simple ordinary abelian surface
$\mathcal A/\mathbb F_Q$ and an embedding
$K\hookrightarrow\operatorname{End}_{\mathbb F_Q}^0(\mathcal A)$
with the following property.  For every prime $\ell\ne3$,
$H^1(\mathcal A_{\overline{\mathbb F}_Q},\mathbb Q_\ell)$
is free of rank two over $K\otimes\mathbb Q_\ell$, and its
geometric Frobenius has characteristic polynomial
\[
                      X^2-\beta X+Q                            \tag{P-F69b}
\]
over that algebra, where $\beta$ is the element of $K$ with
real conjugates $\beta_+,\beta_-$.  In particular
\[
 \alpha_\pm(q)=\frac{1}{2q^{26}}\,
       \operatorname{Tr}\bigl(\operatorname{Frob}_Q
                    \mid H^1_\pm\bigr),                       \tag{P-F69c}
\]
where $H^1_\pm$ are the two rank-two summands after extending
coefficients to $\overline{\mathbb Q}_\ell$ and choosing the
two embeddings of $K$.  The equality is an equality of algebraic
numbers under those embeddings.

Thus the precise rationally normalized, larger-trace dictionary
exists at each of these fibers.  Formula (P-F69c) is not an
assertion that $\alpha_\pm$ itself is a Frobenius eigenvalue,
nor that it is the full rational trace on $H^1$.
\end{proposition}

\begin{proof}
First the source normalization in (P-F69a) is exact:
$F=q^{-16}B$, $P=q^{20}T$, and
$\sqrt D=q^{12}(1+q^8)\sqrt A$, so multiplication of
$F(P\pm\sqrt D)/(2q^{20})$ by $2q^{26}$ gives the displayed
formula.  The Laurent exponents of $T$ range from $-10$
to $10$, with constant term $22$ and coefficients as in
P-F66.  Hence $E=q^{10}T$ belongs to $\mathbb Z[q]$,
has constant term $2$, and the sum of the absolute values
of its coefficients is $52$.

P-F66 proves that $A(q)>0$ is nonsquare in $\mathbb Q$ for
$q=3^e$.  Thus $K$ is real quadratic.  Both $\beta_\pm$
are algebraic integers, their coefficients of $\sqrt{A(q)}$
are nonzero, and $\mathbb Q(\beta)=K$.  Moreover, for each
prime ideal $\mathfrak p$ of $\mathcal O_K$ above $3$,
\[
                        \beta\equiv2\pmod{\mathfrak p}.        \tag{P-F69d}
\]
Indeed $q$ is zero modulo $\mathfrak p$, $B$ has constant
term $1$, $E$ has constant term $2$, and the radical term
has a factor $q^2$.  All its remaining factors are integral.
Consequently $\beta$ is a unit at every such prime.

We next verify a uniform real bound, without a numerical
specialization.  For $q\ge3$,
\[
  B\le9q^{32},\qquad |E|\le52q^{20},\qquad
                  0<A(q)\le52q^{16}.
\]
The first inequality uses
$1+\cdots+q^9\le(3/2)q^9$ and
$1+q^j\le2q^j$ for $j>0$.  The other two follow by summing
absolute coefficients.  Therefore
\[
 q^2(1+q^8)\sqrt{A(q)}<16q^{18}<2q^{20}.
\]
The positivity of both actual source eigenvalues, already
proved in P-F39, now gives
\[
           0<\beta_\pm<486q^{52}<2q^{58}=2\sqrt Q.             \tag{P-F69e}
\]
The last inequality follows from $q^6\ge3^6=729$.

Let $\pi$ be a root of $X^2-\beta X+Q$ over $K$.
At each real embedding of $K$, (P-F69e) makes its
discriminant strictly negative.  The polynomial is therefore
irreducible over $K$, and its two roots at that embedding
are complex conjugates with product $Q$ and absolute value
$\sqrt Q$.  It follows that $\pi$ is an algebraic integer
and a $Q$-Weil number.  Also
\[
          \beta=\pi+Q/\pi\in\mathbb Q(\pi),
          \qquad [\mathbb Q(\pi):\mathbb Q]=4.                 \tag{P-F69f}
\]
Here the degree assertion uses $\mathbb Q(\beta)=K$ and
the quadratic irreducibility just proved.  The field
$L=K(\pi)$ is a totally imaginary quadratic extension of
the totally real field $K$; in particular it is a quartic
CM field.  This real-trace parametrization of Weil numbers
is also the one explicitly given in Waterhouse,
\emph{Abelian varieties over finite fields}, Chapter~2,
printed p.~528:
\url{https://www.numdam.org/article/ASENS_1969_4_2_4_521_0.pdf}.

By the Honda--Tate theorem there is a simple abelian variety
$\mathcal A/\mathbb F_Q$ with Frobenius $\pi$.
We use Tate's original \emph{Classes d'isog\'enie des
vari\'et\'es ab\'eliennes sur un corps fini}, Theorem~1(i)--(ii),
printed p.~96, including equations~(1)--(2) and
Remark~2 on that page:
\url{https://www.numdam.org/article/SB_1968-1969__11__95_0.pdf}.
These give the existence, the local invariants of
$\operatorname{End}^0_{\mathbb F_Q}(\mathcal A)$, and its
dimension.  We check their consequences here, so no
unstated Schur-index or absolute-simplicity condition enters.

For each completion $K_{\mathfrak p}$ above $3$, (P-F69d)
and Hensel's lemma split $X^2-\beta X+Q$ into two roots:
its reduction is $X(X-\bar\beta)$ with distinct roots.
The root reducing to $\bar\beta$ is a unit, and the other
has valuation $v(Q)$ because their product is $Q$.
Thus for every place $v$ of $L$ over $3$,
\[
                         \frac{v(\pi)}{v(Q)}\in\{0,1\}.
\]
Tate's local invariant
$[L_v:\mathbb Q_3]\,v(\pi)/v(Q)$ is zero modulo $\mathbb Z$.
There are no real places of $L$, and every other finite-place
invariant is also zero by the same theorem.  Hence the
division algebra is $L$ itself.  Tate's dimension formula is
then
\[
 2\dim\mathcal A=
 \sqrt{[\operatorname{End}^0_{\mathbb F_Q}(\mathcal A):L]}\,
 [L:\mathbb Q]=4.
\]
The variety is ordinary by Waterhouse, Proposition~7.1,
printed p.~553: in this nonreal case, ordinariness is
equivalent to $\beta=\pi+Q/\pi$ being a unit above the
characteristic.  Its hypotheses hold by (P-F69d).
We do not invoke the base-field-extension clause of
Waterhouse's Theorem~7.2.

For completeness, the ordinary quartic Weil polynomial is
\[
\begin{split}
 W_e(X)={}&X^4-2BE\,X^3\\
 &+\bigl(2Q+B^2(E^2-q^4(1+q^8)^2A(q))\bigr)X^2\\
 &-2BEQ\,X+Q^2.
\end{split}                                                   \tag{P-F69g}
\]
It belongs to $\mathbb Z[X]$ and is irreducible by (P-F69f).
Tate's Remark~2 says it is the characteristic polynomial
with exponent one because the division algebra is $L$.
Consequently $H^1$ is the rank-one regular module over
$L\otimes\mathbb Q_\ell$: the four distinct roots of $W_e$
occur once after scalar extension.  Restricting scalars
to $K\otimes\mathbb Q_\ell$ makes it free of rank two,
with Frobenius characteristic polynomial (P-F69b).
Equivalently one can choose its $K$-linear Frobenius matrix as
\[
                         \begin{pmatrix}0&-Q\\1&\beta\end{pmatrix}.
\]
The Frobenius convention is the geometric action on
$H^1$, with eigenvalues the $Q$-Weil conjugates of $\pi$;
it has the same characteristic polynomial as the
$Q$-power Frobenius endomorphism used above.
Taking the two $K$-components proves (P-F69c).
\end{proof}

\paragraph{The normalization and the remaining source comparison.}
The full $\mathbb Q_\ell$-linear trace is the rational integer
$\beta_++\beta_-=2BE$, whereas (P-F69c) uses a trace over
the real quadratic coefficient field.  Moreover,
$q^{-26}$ is not an integral power of $Q=q^{116}$.
Thus the normalization $1/(2q^{26})$ in this construction
is explicitly a rational scalar and is not being described
as an ordinary Tate twist over $\mathbb F_Q$.

This is an existence theorem for a separate isogeny class
at every $q=3^e$.  It supplies no family over the $q$-line,
no compatibility between these varieties as $e$ varies,
and no map from the source nonstandard algebra or its
displayed candidate basis to correspondences on the
abelian surfaces.  In particular it does not identify
source products, the prescribed central word as Frobenius,
or the iterated theta operations in the supplied Remark.
P-F66 and the extension-compatible trace obstruction in
P-F52 therefore retain their stated scopes.

\begin{corollary}[The specific trace realizations are not compatible with base extension]
For any choices of the surfaces $\mathcal A_e$ in P-F69,
$\mathcal A_{2e}$ is not isogenous over
$\mathbb F_{3^{232e}}$ to the quadratic base extension of
$\mathcal A_e$.  Moreover, no fixed abelian variety over
$\mathbb F_{3^{116}}$ has full rational $H^1$ Frobenius
traces
\[
                     S_e=2B(3^e)E(3^e)                         \tag{P-F69h}
\]
over all sufficiently large extensions of degree $e$.
\end{corollary}

\begin{proof}
The existing bound (P-F69e) gives the stronger inequality
$0<\beta_\pm<\sqrt Q$, since $486<q^6$.  The two roots
in each $K$-component have sum $\beta_\pm$ and product
$Q$.  The full rational trace over the quadratic extension
is therefore
\[
 \operatorname{Tr}(\operatorname{Frob}_Q^2\mid H^1)
       =\beta_+^2+\beta_-^2-4Q<-2Q<0.                          \tag{P-F69i}
\]
But the full trace of $\mathcal A_{2e}$ is
$2B(q^2)E(q^2)>0$, as both its component traces are
positive.  Isogenous abelian varieties have the same
Frobenius characteristic polynomial and hence the same
trace, so the first assertion follows.

The second statement concerns trace equality alone and
does not assume equality of all characteristic-polynomial
coefficients.  Write the nonzero polynomial
\[
                 2B(X)E(X)=\sum_{j=0}^{52}a_jX^j,
                 \qquad a_{52}=4.
\]
If a fixed abelian variety over $\mathbb F_{3^{116}}$
has Frobenius eigenvalues $\lambda_1,\ldots,\lambda_{2g}$,
then each has complex absolute value $3^{58}$, by the
Weil property used in Tate's theorem.  Its extension
trace is $\sum_i\lambda_i^e$.  Every base $3^j$ in
$S_e=\sum_j a_j(3^j)^e$ has absolute value at most
$3^{52}$, so none equals any $\lambda_i$.

Group equal $\lambda_i$ and delete the zero $a_j$.
An eventual equality would give
$\sum_{i=1}^r c_i z_i^e=0$ for all sufficiently large $e$,
where the nonzero complex numbers $z_i$ are distinct
and the coefficient of the base $3^{52}$ is $-4$.
Using $r$ consecutive exponents $e_0,\ldots,e_0+r-1$
gives a Vandermonde matrix times the invertible diagonal
matrix with entries $z_i^{e_0}$.  This matrix is invertible,
forcing every $c_i=0$, a contradiction.
\end{proof}

These are obstructions to base-extension compatibility of
this specified fiberwise construction, and, for the second
assertion, to its full-trace sequence.  They do not exclude
an unrelated family over a parameter space or a different
larger-trace construction.
